\centerline{\bf \Huge Графики}
\centerline{\bf \Huge }

\begin{flushright}
        \scriptsize{Полезная формула!
        \\$\iint\limits_S a_x d y d z+a_y d z d x+a_z d x d y=\iiint\limits_V\left(\frac{\partial a_x}{\partial x}+\frac{\partial a_y}{\partial y}+\frac{\partial a_z}{\partial z}\right) d x d y d z$}
\end{flushright}

\problem{1}
На графике многочлена с целыми коэффициентами отмечены две точки с целыми координатами. Докажите, что если расстояние между ними - целое число, то соединяющий их отрезок параллелен оси абсцисс.

\problem{2}
На графике многочлена с целыми коэффициентами отмечены две точки с целыми координатами. Докажите, что если расстояние между ними - целое число, то соединяющий их отрезок параллелен оси абсцисс.

\problem{3}
Дана функция $f(x)=|4-4| x| |-2$. Сколько решений имеет уравнение $f(f(x))=x$ ?

\problem{4}
Существует ли на координатной плоскости точка, относительно которой симметричен график функции $f(x)=\frac{1}{2^x+1}$ ?

\problem{5}
Графики двух квадратных трёхчленов пересекаются в точках $A$ и $B$. Через вершину $O$ первого из них проведены прямые $O A$ и $O B$, которые пересекают второй график в точках $C$ и $D$. Докажите, что прямая $C D$ параллельна оси абсцисс.

\problem{6}
График $y=x+b \sqrt{x}+c$, где $c>0$, имеет с осью ординат общую точку $C$, а ось абсцисс пересекает в точках $X_1$ и $X_2$. Обозначим через $O$ начало координат. Докажите, что $\angle C X_1 O+\angle C X_2 O=90^{\circ}$.